\chapter{What Makes a Reduced Model Adequate?}
\index{model adequacy|textbf}
\index{reduced model|textbf}

\label{ch:operator-reduction-motivation}

\chapterstatus{conceptual and illustrative framework. The central hypothesis,
continuum criteria, and intellectual contributions are proposed work; the
diagnostic examples are not research evidence or validated reductions of
silicon.}

\section{Project motivation and significance}

This section provides a condensed review of quantum mechanics, solid-state
physics, and semiconductor physics to elucidate how the present research
program grew from a pedagogical question.\footnote{For quantum mechanics, see
Shankar~\cite{shankar1994}. For solid-state physics, see
Kittel~\cite{kittel2004} and Simon~\cite{simon2013}. Sze and
Lee~\cite{szeLee2012} provide a gentle device-oriented introduction to
semiconductor physics, while Yu and Cardona~\cite{yuCardona2010} provide a
graduate-level treatment.} The review is not intended to replace a textbook
treatment. It establishes the operator language needed to connect
the familiar microscopic Hamiltonian to a controlled reduced model.

The starting point is an operator equation, not a matrix or a list of energy
eigenvalues.

\begin{definition}[Linear operator]
\index{operator!definition|textbf}
\index{domain!operator|textbf}

Let $\mathcal H$ be a Hilbert space of quantum states. A linear operator
$\hat A$ is a map
\begin{equation}
    \hat A:\mathcal D(\hat A)\subseteq\mathcal H\longrightarrow\mathcal H,
\end{equation}
with a domain $\mathcal D(\hat A)$ on which its action is defined.
\end{definition}

This definition keeps the abstract operator distinct from any matrix used to
represent it. The domain is part of the operator: writing
$\hat A\lvert\psi\rangle$ is meaningful only when
$\lvert\psi\rangle\in\mathcal D(\hat A)$. That distinction becomes important
when differential Hamiltonians are replaced by finite matrices or restricted
to retained subspaces.

\begin{definition}[Dirac notation for states and operators]
\index{Dirac notation|textbf}
\index{quantum state!ket notation|textbf}
\index{expectation value|textbf}

A quantum state is represented by a ket
$\lvert\psi\rangle\in\mathcal H$. Its corresponding bra
$\langle\psi\rvert$ is the linear functional associated with the ket by the
Hilbert-space inner product. For
$\lvert\phi\rangle,\lvert\psi\rangle\in\mathcal H$, their inner product and the
norm of $\lvert\psi\rangle$ are written
\begin{equation}
    \langle\phi\vert\psi\rangle,
    \qquad
    \lVert\psi\rVert
    =
    \sqrt{\langle\psi\vert\psi\rangle}.
\end{equation}
A normalized state satisfies
$\langle\psi\vert\psi\rangle=1$. If
$\lvert\psi\rangle\in\mathcal D(\hat A)$, the ket
$\hat A\lvert\psi\rangle$ denotes the state obtained by applying $\hat A$.
For a normalized state
$\lvert\psi\rangle\in\mathcal D(\hat A)$ and a self-adjoint observable
$\hat A$, the expectation value is
\begin{equation}
    \langle\hat A\rangle_{\psi}
    =
    \langle\psi\rvert\hat A\lvert\psi\rangle.
\end{equation}
\end{definition}

Dirac notation separates a state from any particular coordinate description of
that state. If $\{\lvert\phi_a\rangle\}$ is an orthonormal basis, then
\begin{equation}
    \lvert\psi\rangle
    =
    \sum_a c_a\lvert\phi_a\rangle,
    \qquad
    c_a=\langle\phi_a\vert\psi\rangle.
\end{equation}
The ket is the abstract state; the coefficients $c_a$ are its coordinates in
the selected basis. Changing the basis changes the coordinates without, by
itself, changing the state. This elementary distinction becomes essential when
Bloch, Wannier, tight-binding, and continuum representations are compared.

\begin{definition}[Hamiltonian and additive decomposition]
\index{Hamiltonian!definition|textbf}
\index{Hamiltonian!additive decomposition|textbf}

A Hamiltonian $\hat H$ is the operator that generates the dynamics of the
selected physical model. For the elementary nonrelativistic decomposition used
here,
\begin{equation}
    \hat H=\hat T+\hat V,
\end{equation}
where $\hat T$ is the kinetic-energy operator and $\hat V$ is the potential or
interaction operator. For every state $\lvert\psi\rangle$ in their common
domain,
\begin{equation}
    \hat H\lvert\psi\rangle
    =
    \bigl(\hat T+\hat V\bigr)\lvert\psi\rangle.
\end{equation}
\end{definition}

This is an equality of operator actions on states. It is not yet a matrix
equality, because no finite subspace or basis has been selected.

\begin{definition}[Orthogonal projection]
\index{projection!orthogonal, definition|textbf}
\index{retained subspace!orthogonal projector|textbf}

Let $\mathcal S\subseteq\mathcal H$ be a closed subspace. Its orthogonal
projector $\hat P_{\mathcal S}$ is the operator satisfying
\begin{equation}
\begin{aligned}
    \text{(i)}\quad
    &\hat P_{\mathcal S}^{2}=\hat P_{\mathcal S},
    &&\text{(idempotence)},\\
    \text{(ii)}\quad
    &\hat P_{\mathcal S}^{\dagger}=\hat P_{\mathcal S},
    &&\text{(self-adjointness)},\\
    \text{(iii)}\quad
    &\operatorname{Ran}\hat P_{\mathcal S}=\mathcal S,
    &&\text{(range condition)}.
\end{aligned}
\end{equation}
\end{definition}

For a finite orthonormal basis
$\{\lvert\phi_a\rangle\}_{a=1}^{N}$ of $\mathcal S$, the projector has the
representation
\begin{equation}
    \hat P_{\mathcal S}
    =
    \sum_{a=1}^{N}
    \lvert\phi_a\rangle\langle\phi_a\rvert.
\end{equation}
Each term retains the component of a state along
$\lvert\phi_a\rangle$, while the sum removes every component orthogonal to
$\mathcal S$.

The compressed Hamiltonian is
$\hat P_{\mathcal S}\hat H\hat P_{\mathcal S}$ restricted to $\mathcal S$.
Consistently compressing the decomposition gives
\begin{equation}
    \hat P_{\mathcal S}\hat H\hat P_{\mathcal S}
    =
    \hat P_{\mathcal S}\hat T\hat P_{\mathcal S}
    +
    \hat P_{\mathcal S}\hat V\hat P_{\mathcal S}.
\end{equation}
Projection may change the mathematical form of the terms: even when $\hat V$
is local in the original coordinate representation, its compression need not
be representable by multiplication by a scalar function on the retained
space.

The projector augmented-wave (PAW) method provides
\index{PAW method@projector augmented-wave (PAW) method}
\index{transformation!PAW}
\index{overlap metric!PAW}
 an instructive contrast to
this orthogonal compression. In PAW, a linear transformation
$\hat{\mathcal T}$ maps a smooth auxiliary Kohn--Sham state to an all-electron
state,
\begin{equation}
    \lvert\psi_n\rangle
    =
    \hat{\mathcal T}\lvert\widetilde\psi_n\rangle.
\end{equation}
The auxiliary state consequently satisfies the generalized eigenproblem
\begin{equation}
    \hat{\mathcal T}^{\dagger}\hat H\hat{\mathcal T}
    \lvert\widetilde\psi_n\rangle
    =
    \varepsilon_n
    \hat{\mathcal T}^{\dagger}\hat{\mathcal T}
    \lvert\widetilde\psi_n\rangle.
\end{equation}
rather than an ordinary orthonormal eigenproblem~\cite{rostgaard2009}. The
localized PAW projector functions recover expansion coefficients inside
augmentation regions; despite their name, they are not the global orthogonal
projector $\hat P_{\mathcal S}$ defined above. Rostgaard's tutorial also shows
how the transformation must be carried into expectation values and Wannier
matrix elements. PAW therefore illustrates why a state transformation, an
induced overlap metric, and an orthogonal subspace projection must not be
silently identified.

The active silicon parent uses approved norm-conserving pseudopotentials rather
than PAW datasets. This comparison is explanatory and does not alter that
scientific setting.

\readingnote{Read Bl\"ochl's original PAW paper~\cite{blochl1994paw} before
using PAW transformations, projector functions, or augmentation corrections in
a scientific claim. Rostgaard~\cite{rostgaard2009} is a useful self-contained
tutorial, but the original paper is the primary source for the method.}

\begin{definition}[Position-space representation]
\index{position-space representation|textbf}
\index{coordinate representation|textbf}

Let $\langle\mathbf r\rvert$ be the generalized position bra and define the
wavefunction by
$\psi(\mathbf r)=\langle\mathbf r\vert\psi\rangle$. The position-space
representation of the abstract operator equation is
\begin{equation}
    \langle\mathbf r\rvert\hat H\lvert\psi\rangle
    =
    \langle\mathbf r\rvert\hat T\lvert\psi\rangle
    +
    \langle\mathbf r\rvert\hat V\lvert\psi\rangle.
\end{equation}
For the nonrelativistic kinetic operator and a local scalar potential, this
becomes
\begin{equation}
    \langle\mathbf r\rvert\hat H\lvert\psi\rangle
    =
    -\frac{\hbar^{2}}{2m}\nabla^{2}\psi(\mathbf r)
    +V(\mathbf r)\psi(\mathbf r).
\end{equation}
\end{definition}

The generalized coordinate bra $\langle\mathbf r\rvert$ is not itself the
finite-rank projector used above. Coordinate representation, subspace
projection, basis selection, and model reduction are distinct operations.
Maintaining those distinctions is the foundation of this monograph.

Semiconductor physics is commonly taught through two quantum-mechanical
descriptions that are physically related but mathematically disconnected.
Students first solve a microscopic periodic Hamiltonian,
\begin{equation}
    \hat H_{\mathrm{micro}}
    =
    -\frac{\hbar^2}{2m_0}\nabla^2
    +
    V_{\mathrm{per}}(\mathbf r),
\end{equation}
and obtain Bloch states and energy bands. They are later introduced to
effective-mass theory through a continuum envelope Hamiltonian. Written
relative to the selected band-edge energy, it is
\begin{equation}
    \hat H_{\mathrm{EMT}}
    =
    -\frac{\hbar^2}{2}
    \nabla\mathbin{\cdot}\mathbf m^{*-1}\nabla
    +
    V_{\mathrm{ext}}(\mathbf r).
\end{equation}

For a smooth, nondegenerate band near an extremum $\mathbf k_0$, the usual
connection expands the dispersion as
\begin{equation}
    E_n(\mathbf k_0+\mathbf q)
    =
    E_n(\mathbf k_0)
    +
    \frac{\hbar^2}{2}
    \mathbf q^{\mathsf T}\mathbf m_n^{*-1}\mathbf q
    +
    O\!\left(\lVert\mathbf q\rVert^3\right).
\end{equation}
After subtracting the constant $E_n(\mathbf k_0)\hat I$ and replacing
$\mathbf q$ by $-i\nabla$, this argument explains the form of the effective
kinetic operator. It does not construct the envelope-function state space,
define the map from microscopic Bloch states to continuum states, or quantify
the approximations introduced by band truncation, gauge selection, valley
decomposition, and spatial coarse-graining.

The Luttinger--Kohn derivation, its assumptions, its multivalley and degenerate-
band extensions, and its relation to the present first-principles reduction
problem are developed in Appendix~\ref{app:envelope-theory}. The detailed
summary and remaining source-review note are kept there rather than duplicated
in this chapter~\cite{luttingerKohn1955}.

Effective-mass theory is therefore often presented as a second quantum
mechanics rather than as a controlled reduction of the first. The
effective-mass tensor is obtained from the curvature of a Bloch band, but the
corresponding state-space reduction is usually left implicit.

This project began from a practical attempt to close that pedagogical gap. The
initial objective was to construct an explicit periodic potential that students
could discretize and diagonalize, yet which remained sufficiently expressive
to produce an indirect band gap. Such a model would provide an executable
bridge between elementary periodic quantum mechanics and the band structure of
an indirect-gap semiconductor.

The most direct strategy appears elementary: calculate Kohn--Sham
density-functional-theory Bloch eigenpairs, reconstruct a retained Hamiltonian,
subtract the known kinetic operator, and interpret the remainder as an
effective periodic potential.

Let $\hat H_{\mathrm{KS}}(\mathbf k)$ be a Kohn--Sham Bloch-fiber operator
with cell-periodic eigenstates $\lvert u_{n\mathbf k}\rangle$. For a retained
set of bands $\mathcal B$, define
\begin{equation}
    \hat P(\mathbf k)
    =
    \sum_{n\in\mathcal B}
    \lvert u_{n\mathbf k}\rangle
    \langle u_{n\mathbf k}\rvert
\end{equation}
and
\begin{equation}
    \hat H_{\mathrm{KS}}^{(P)}(\mathbf k)
    =
    \left.
        \hat P(\mathbf k)
        \hat H_{\mathrm{KS}}(\mathbf k)
        \hat P(\mathbf k)
    \right|_{\mathcal H^{(P)}(\mathbf k)}
    =
    \sum_{n\in\mathcal B}
    \varepsilon_n(\mathbf k)
    \lvert u_{n\mathbf k}\rangle
    \langle u_{n\mathbf k}\rvert,
\end{equation}
where
\begin{equation}
    \mathcal H^{(P)}(\mathbf k)
    =
    \operatorname{Ran}\hat P(\mathbf k).
\end{equation}
For entangled bands, constructing a smooth retained family may require
projection or disentanglement rather than selection by a single global energy
cutoff.

The apparently obvious construction is
\begin{equation}
    \hat V_{\mathrm{eff}}^{(P)}(\mathbf k)
    \stackrel{?}{=}
    \hat H_{\mathrm{KS}}^{(P)}(\mathbf k)
    -
    \hat T^{(P)}(\mathbf k),
\end{equation}
where
\begin{equation}
    \hat T^{(P)}(\mathbf k)
    =
    \left.
        \hat P(\mathbf k)
        \hat T(\mathbf k)
        \hat P(\mathbf k)
    \right|_{\mathcal H^{(P)}(\mathbf k)}.
\end{equation}
This subtraction is mathematically valid on the retained fiber, but its result
is a projected operator, not necessarily a unique local scalar potential.
Projection may induce nonlocality, while nonlocal pseudopotentials, orbital
structure, spin--orbit interactions, and band truncation may introduce
components that cannot be represented by a simple function $V(\mathbf r)$.
Bloch gauge freedom further changes the matrix coordinates of the retained
operator without changing its abstract action or spectrum.

\citationtodo{medium priority}{Check the entangled-band disentanglement and
momentum-dependent Wannier-gauge statements against Souza--Marzari--Vanderbilt
(2001) and Marzari et al. (2012), bibliography keys \texttt{souza2001} and
\texttt{marzari2012} (already present).}

The desired teaching potential must therefore be obtained by reducing the
operator residual into an explicitly chosen model class:
\begin{equation}
    \hat V_{\mathrm{teach}}
    =
    \operatorname*{arg\,min}_{\hat V\in\mathcal V_{\mathrm{adm}}}
    \sum_{\mathbf k\in\mathcal K}
    \left\|
        \hat H_{\mathrm{KS}}^{(P)}(\mathbf k)
        -
        \hat T^{(P)}(\mathbf k)
        -
        \left.
            \hat P(\mathbf k)
            \hat V(\mathbf k)
            \hat P(\mathbf k)
        \right|_{\mathcal H^{(P)}(\mathbf k)}
    \right\|_*^2.
\end{equation}
Here, $\mathcal K$ is the sampled reciprocal-space region and
$\mathcal V_{\mathrm{adm}}$ may represent local periodic potentials, truncated
Fourier expansions, finite real-space stencils, or constrained tight-binding
operator classes. The norm $\lVert\cdot\rVert_*$ and any reciprocal-space
weights must be chosen to measure operator error consistently with the physical
observables and continuum limit of interest. This is a proposed model-class
problem, not a presently accepted fitting criterion.

The finite example developed next separates exact operator retention from
model-class approximation and shows why all components of an additive
Hamiltonian must be reduced through the same map. The later sections then
return to pristine--doped comparison, continuum admissibility, and the broader
research program.

\section{From spectral data to an effective Hamiltonian}
\index{spectral reconstruction}
\index{effective Hamiltonian}


The construction in this section is a pedagogical discrete-spectrum case. It
is not the literal subspace construction used for a periodic crystal, where
Bloch-fiber representations and possibly projection or disentanglement enter.

The genesis of this problem begins with the spectral problem for a quantum
Hamiltonian $\hat H$,
\begin{equation}
    \hat H\lvert\psi_i\rangle
    =
    E_i\lvert\psi_i\rangle,
\end{equation}
where $(E_i,\lvert\psi_i\rangle)$ is the $i$th eigenpair. Assuming that the
relevant spectrum is discrete and that the eigenstates form a complete
orthonormal basis of the Hilbert space $\mathcal H$, the Hamiltonian admits the
spectral decomposition
\begin{equation}
    \hat H
    =
    \sum_i
    E_i\lvert\psi_i\rangle\langle\psi_i\rvert.
\end{equation}

Suppose that the physical questions of interest involve only states below an
energy cutoff $E_{\mathrm{cut}}$. We define the retained low-energy subspace
\begin{equation}
    \mathcal H^{(P)}
    =
    \operatorname{span}
    \left\{
        \lvert\psi_i\rangle
        \;\middle|\;
        E_i\leq E_{\mathrm{cut}}
    \right\},
\end{equation}
with orthogonal projector
\begin{equation}
    \hat P
    =
    \sum_{E_i\leq E_{\mathrm{cut}}}
    \lvert\psi_i\rangle\langle\psi_i\rvert.
\end{equation}

The corresponding retained Hamiltonian is
\begin{equation}
    \hat H^{(P)}
    =
    \left.
        \hat P\hat H\hat P
    \right|_{\mathcal H^{(P)}}
    =
    \sum_{E_i\leq E_{\mathrm{cut}}}
    E_i\lvert\psi_i\rangle\langle\psi_i\rvert.
\end{equation}
It exactly reproduces every retained eigenpair,
\begin{equation}
    \hat H^{(P)}\lvert\psi_i\rangle
    =
    E_i\lvert\psi_i\rangle,
    \qquad
    E_i\leq E_{\mathrm{cut}},
\end{equation}
but it does not specify the action of $\hat H$ outside
$\mathcal H^{(P)}$. It is therefore an exact restriction of the parent
Hamiltonian, not yet an approximate effective model.

We reserve the notation $\hat H_{\mathrm{eff}}(\boldsymbol\theta)$ for an
operator belonging to a prescribed effective-model class with parameters
$\boldsymbol\theta$. The distinction is important:
\begin{equation}
    \underbrace{\hat H^{(P)}}_{\text{exact retained operator}}
    \qquad\text{and}\qquad
    \underbrace{\hat H_{\mathrm{eff}}(\boldsymbol\theta)}_{
        \text{approximate model on an identified retained space}}
\end{equation}
are not the same construction.

\section{Controlled diagnostic examples}
\index{diagnostic example}


Simple finite-dimensional examples expose the distinction between matching a
spectrum and identifying an operator. Appendix~\ref{app:particle-box} develops
the particle-in-a-box residual and its identifiability limits in detail, while
Appendix~\ref{app:harmonic-oscillator} compares two finite harmonic-oscillator
constructions. These examples are analytical and numerical diagnostics; they
are not evidence for the silicon reduction program.

\section{From the particle in a box to electronic-structure reduction}

The particle-in-a-box example is elementary, but the obstruction it reveals is
general. In a common parent state space, the decomposition
\begin{equation}
    \mathbf H=\mathbf T+\mathbf V
\end{equation}
is meaningful because all three matrices act on the same discrete state space
and are represented in the same basis. Reduction preserves this additive
decomposition only when every term is reduced through the same map:
\begin{equation}
    \mathcal R_{P_M}(\mathbf H-\mathbf T)
    =
    \mathcal R_{P_M}(\mathbf H)
    -
    \mathcal R_{P_M}(\mathbf T).
\end{equation}
In general,
\begin{equation}
    \mathcal R_{P_M}(\mathbf H-\mathbf T)
    \neq
    \mathcal R_{P_M}(\mathbf H)-\mathbf T.
\end{equation}

The distinction becomes more consequential in electronic-structure
calculations. A Kohn--Sham Hamiltonian, a Wannier Hamiltonian, a tight-binding
Hamiltonian, and an effective-mass Hamiltonian generally do not begin as
matrices acting on an identical numerical state space. They may differ in
dimension, basis, gauge, spatial resolution, boundary conditions, geometry,
units, and energy reference. Their matrices therefore cannot be physically
compared merely because they are all called Hamiltonians.

The central problem of this work is consequently not the mechanical
subtraction of two matrices. It is the construction and validation of the
state-space identifications that make such a subtraction mathematically
defined and physically interpretable.

\section{From periodic potentials to impurity operators}
\index{impurity operator!motivation}
\index{periodic potential}


The same structural problem becomes more consequential when the model is
extended from bulk bands to semiconductor impurities. Given pristine and doped
Hamiltonians, the natural proposal is
\begin{equation}
    \hat V_{\mathrm{imp}}
    \stackrel{?}{=}
    \hat H_d-\hat H_b.
\end{equation}
This expression is not yet an operator identity because independently
constructed pristine and doped calculations may use different retained
subspaces, Bloch or Wannier gauges, basis orderings, geometries, spin
structures, energy references, and numerical dimensions. Their matrices cannot
be interpreted as representations in common physical coordinates without an
explicit identification.

Let $s\in\{b,d\}$ denote matched bulk and doped branches, with retained
operators
\begin{equation}
    \hat H_s^{(P)}
    =
    \left.
        \hat P_s\hat H_s\hat P_s
    \right|_{\mathcal H_s^{(P)}},
    \qquad
    \mathcal H_s^{(P)}
    =
    \operatorname{Ran}\hat P_s.
\end{equation}
If the retained ranks agree, a meaningful comparison requires an
identification
\begin{equation}
    \hat U_d:
    \mathcal H_b^{(P)}
    \longrightarrow
    \mathcal H_d^{(P)}.
\end{equation}
Let $c_b$ and $c_d$ be the scalar reference energies selected by the declared
alignment procedure, and define
\begin{equation}
    \overline{\hat H}_s^{(P)}
    =
    \hat H_s^{(P)}-c_s\hat I_s.
\end{equation}
For a unitary identification $\hat U_d$, the aligned impurity operator on the
bulk retained space is
\begin{equation}
    \Delta\hat H_d^{(P)}
    =
    \hat U_d^\dagger
    \overline{\hat H}_d^{(P)}
    \hat U_d
    -
    \overline{\hat H}_b^{(P)}.
\end{equation}
Unequal retained ranks require a separately defined partial identification or
further model reduction; they do not admit this unitary formula unchanged.

Equal-dimensional spaces admit many abstract unitary maps. Most are not
physically admissible: an arbitrary unitary can scramble localization, orbital
character, crystal momentum, valley content, symmetry, and spin while leaving
the spectrum unchanged. Gauge alignment is therefore necessary to define the
subtraction, but it is not sufficient to select a physical impurity operator.
The identification must also respect the structures needed by the intended
reduced theory.

\section{Central hypothesis}
\index{central hypothesis}
\index{model class!admissible}


The central hypothesis proposed by this project is:
\begin{quote}
\itshape
The intended continuum limit is not merely a downstream approximation to an
already constructed atomistic model. It constrains, from the beginning, the
admissible retained spaces, gauges, identification maps, and reduced operator
classes.
\end{quote}

For a semiconductor representation to support an effective-mass
interpretation, its atomistic-to-continuum identification must preserve or
control the structures used by the continuum theory. These include
\begin{itemize}
    \item real-space localization and lattice translation;
    \item Bloch wave vector and valley correspondence;
    \item orbital, symmetry, and spin content;
    \item smoothness across the relevant reciprocal-space region;
    \item common bulk behavior away from an impurity;
    \item boundary conditions and normalization; and
    \item separation of slowly varying envelopes from lattice-periodic
          components.
\end{itemize}

Let
\begin{equation}
    \epsilon=\frac{a}{L}
\end{equation}
denote the ratio of lattice scale $a$ to an envelope or impurity scale $L$.
For a specified family of retained atomistic operators
$\hat H_\epsilon^{(P)}$ and appropriately normalized identification maps
\begin{equation}
    \hat J_\epsilon:
    \mathcal H_\epsilon^{(P)}
    \longrightarrow
    \mathcal H_{\mathrm{EMT}},
\end{equation}
one possible continuum criterion is norm-resolvent convergence,
\begin{equation}
    \left\|
        \hat J_\epsilon
        \left(\hat H_\epsilon^{(P)}-z\right)^{-1}
        \hat J_\epsilon^\dagger
        -
        \left(\hat H_{\mathrm{EMT}}-z\right)^{-1}
    \right\|
    \longrightarrow 0,
    \qquad
    \epsilon\longrightarrow 0,
\end{equation}
for $z$ in an identified common resolvent set. This is a candidate analytical
form, not an established result of the project. A later specification may
instead require a weaker, explicitly defined combination of spectral,
subspace, operator, and observable convergence criteria.

\citationtodo{low priority}{Check the mathematical meaning and implications of
norm-resolvent convergence with an operator-theory source. Candidate: Teschl
(2014), bibliography key \texttt{teschl2014} (already present). The project's
choice of this criterion remains explicitly provisional.}

The continuum limit thereby becomes an upstream admissibility condition. A
gauge, basis, identification, or reduced operator class is acceptable only if
it preserves, or approximates with controlled error, the microscopic structures
needed to recover the declared envelope theory.

\section{Relation to Hilbert's sixth problem}
\index{Hilbert's sixth problem}


This research is situated within the broad mathematical tradition associated
with Hilbert's sixth problem: the rigorous formulation of physical theories and
the justification of limiting relations between microscopic and continuum
descriptions. It does not claim to solve Hilbert's sixth problem in its general
form. Instead, it addresses a bounded and experimentally consequential
instance: the reduction of microscopic electronic-structure operators to
continuum envelope models for indirect-gap semiconductors and their impurities.

\citationtodo{medium priority}{Check this characterization against Hilbert's
published problem statement rather than an uncited historical paraphrase.
Candidate: Hilbert, ``Mathematical Problems'' (1902), provisional key
\texttt{hilbert1902}; the candidate record is present for this editorial note.}

The relevant limiting process is not only a passage between different length
scales. It is a passage between different state spaces, operator algebras,
coordinate systems, and model classes:
\begin{equation}
    \begin{aligned}
        \text{KS--DFT}
        &\longrightarrow \text{retained Bloch operator}\\
        &\longrightarrow \text{localized atomistic operator}\\
        &\longrightarrow \text{effective-mass continuum operator}.
    \end{aligned}
\end{equation}
Each arrow requires an explicit map, a statement of preserved structure, and an
error criterion. Without these elements, agreement between selected band
energies does not establish that two models represent the same reduced physics.

The effective-mass limit is therefore both an endpoint and a constraint. The
program asks not only whether an atomistic calculation can reproduce a
continuum observable, but which atomistic representations and reductions are
capable of supporting the continuum theory at all.

\section{Intellectual contribution}

The proposed program aims to develop a gauge-equivariant,
continuum-constrained framework for semiconductor operator reduction. Its
intended contributions are
\begin{enumerate}
    \item reconstruction of retained operators from Kohn--Sham Bloch
          eigenpairs;
    \item separation of operator reconstruction from model-class
          approximation;
    \item quantitative projection onto interpretable periodic-potential,
          stencil, and tight-binding classes;
    \item explicit alignment of pristine and doped retained state spaces;
    \item construction of impurity operators whose abstract action is invariant
          under compatible coordinate changes after alignment;
    \item formulation of continuum-compatibility criteria using operator,
          subspace, spectral, and observable errors; and
    \item an executable pedagogical chain connecting periodic potentials,
          indirect band structure, impurity operators, and effective-mass
          theory.
\end{enumerate}

The conceptual distinction is concise:
\begin{quote}
\itshape
Gauge alignment makes an operator comparison mathematically defined;
continuum compatibility makes the resulting reduced model physically
admissible.
\end{quote}

What began as an attempt to construct a teachable periodic potential therefore
exposes a broader mathematical problem. Obtaining the potential requires
reconstructing an operator rather than merely fitting a spectrum. Introducing
an impurity requires comparing independently constructed operators. Comparing
those operators requires identifying their state spaces and gauges. Selecting
physically admissible identifications and reduced models requires a declared
continuum limit.

This chain provides both the scientific rationale and the educational
significance of the proposed research. It converts the normally implicit
transition from Bloch theory to effective-mass theory into an explicit,
testable, and executable sequence of operator reductions while keeping every
approximation and evidence claim explicit. Chapter~\ref{ch:physical-problem}
next specializes this framework to the silicon host and its substitutional
phosphorus and boron branches.
